\begin{abstract}
Per-sample gradient ``consensus'' filtering---eigendecomposing the $B \times
B$ inter-sample gradient similarity matrix and keeping only directions above
a Marchenko--Pastur-style threshold---has shown strong label-noise robustness
in classification. We ask whether this mechanism reduces the effect of noise
when training on large, low signal-to-noise financial data. Under a fully
pre-registered protocol on the Numerai v5 dataset (era-purged
tuning/verdict split, frozen thresholds, moving-block-bootstrap inference,
compute-matched tuning), the answer is negative: we find \emph{no evidence of
benefit under the affordable tuning budget}, and at the pre-registered
operating point the filter significantly degraded out-of-sample performance
on both architectures tested (MLP: $-0.00527$ mean per-era correlation, 95\%
CI $[-0.00886, -0.00181]$; GRU: $-0.00484$, CI $[-0.00758, -0.00211]$,
against a $+0.0064$ baseline edge). Two mechanistic contributions explain why
the label-noise story does not transfer. First, we prove that for
scalar-output MSE the filter's engagement eigenspectrum is exactly
target-independent: per-sample gradients reduce to $\pm$(residual sign)
$\times$ a target-independent Jacobian direction, so ``consensus'' measures
feature/factor structure, not signal (verified in fp64 to $8.9 \times
10^{-16}$; independently re-verified). Second, matched controls show
MP-eigenselection is statistically indistinguishable from a norm- and
$k$-matched \emph{random}-subspace projection on both architectures. An
independent audit re-executed the main experiment with fresh seeds and an
independently constructed evaluation shard and reproduced the result. Claims
are scoped to a subsampled, low-edge regime.
\end{abstract}
